\documentclass[../report.tex]{subfiles}

\begin{document}

\subsection{Windows Scanning and Buffer Overflow}

Something that cost time unnecessarily was a lack of attention to detail during the first appointment.
Executing the script kept failing to write the "bad character" candidate string to memory, no matter what characters were passed. I failed to notice for a while that the first part of the string, the \enquote{\textit{AUTH}} portion, was missing a whitespace (see \cref{miss_space}). From that point onward, I was more vigilant.
\begin{figure}[!h]
\includegraphics[width=\linewidth]{buffer_overflow/reflection/missing_whitespace.png}
\caption{Missing whitespace in AUTH string}
\label{miss_space}
\end{figure}

Another misstep was the assumption, that the IP address on the eth1 network interface should be used in payload generation (see \cref{miss_ip}).
\begin{figure}[!h]
\includegraphics[width=\linewidth]{buffer_overflow/reflection/ip}
\caption{Wrong IP address for the attack}
\label{miss_ip}
\end{figure}

Thinking back to the networking lecture \enquote{Rechnernetze}, I realized all addresses are located inside private address ranges. That made me consider them actively, which lead to the realization that the address the computer was assigned inside the ICS laboratory network under the virtual \textbf{tun0} interface is the address that was necessary for the reverse shell connection to work.

Starting the VPN at the beginning of each session was an afterthought, which lead to me not realizing it impacts how the machine needs to be addressed for the attack to work. 
This realization took a while to manifest, and made it clear to me that I should question my assumptions, and make sure that the conclusions I reach are thought through instead of assumed without critical questioning.

Using the ImmunityDebugger gave me a little trouble at first.
A lot of time was spent during the first session learning the tooling and accustoming myself to the environment.

Exploiting the Buffer Overflow required understanding of the stack and how operations work on the machine level, which I lacked from prior studies. I had grasped the concept from explanations during the lecture, but did not fully understand how to apply the concept in practice.
That became clear during work on the assignment, which helped me gain a better understanding of the processes happening close to the hardware.

Writing the report was very time intensive, requiring more time than I had anticipated when starting out.
In addition, the level of detail required was only apparent once I had started the writing process. When I was documenting my approach, I didn't take enough care to represent the steps I took with readable, appropriately formatted screenshots. I had to go back and repeat some steps to take satisfactory screengrabs, only gaining usable material during the second session. At that point I had gone over the material again and reviewed the steps i took in detail. I was therefore able to repeat the steps quickly with intent. A deeper understanding of the subject matter gave me with the ability to anticipate what information was relevant, which made it easier to take the steps in an easily documentable way.

Moving into the next assignment, I knew what to expect in terms of documentation of my actions. I practiced the discussed concepts on my private system, to prepare myself for the scheduled appointments in which we could work on the assignment. I had learned that preparation and familiarity with the tooling and the subject matter enable me to focus on executing the necessary steps while being able to document them. If I don't know what I am doing or why I am doing it, my ability to document and communicate my actions diminishes drastically. Hence I did my utmost to be prepared for the following assignments. 

\end{document}